\section{Related work}
\label{sec:related}

Language-model agents already plan and run experiments in chemistry \citep{boiko2023coscientist}. In biology they
design genetic perturbation screens, tested against data from screens already completed \citep{roohani2024biodiscovery}.
\citet{biodesignbench2026} test agents on protein design and trace their shortfall to how thoroughly they check
candidates. The agents pick suitable tools, score each candidate on only a few measures and stop exploring early. In
those tasks chance enters through design tools that sample new candidates. In ours it enters through measurement, so
one candidate measured twice gives two different numbers. \citet{wang2026agentnoise} add noise to what users say and
what tools return in existing agent benchmarks, and measure how well agents hold up.

\citet{krishnamurthy2024explore} placed GPT-3.5, GPT-4 and Llama2 in multi-armed bandits, where an agent repeatedly
picks one of several options with random payoffs. The prompt held the task and the history of past picks, either raw
or summarised. Only GPT-4, reasoning step by step with a summarised history, explored well. Our ledger gives the same
kind of summary, and in our study it changed re-testing little (Section~\ref{sec:why}). Their payoffs are coin flips,
so the noise level is tied to the payoff and cannot be set on its own.

LLAMBO puts a language model inside Bayesian optimisation, to propose candidates and predict their scores
\citep{liu2024llambo}. \citet{gupta2025bayesopt} tested whether such a model uses the results it gets: shuffling the
measured outcomes on genetic perturbation tasks left performance essentially unchanged. On feedback from a real CRISPR
screen, \citet{wainrib2026labloop} found that shuffling removed the gain. In both studies the results come from
completed screens, and re-measuring a gene is not offered as a use of the budget.

Experimental design has long studied how to split a budget between new conditions and repeats.
\citet{binois2019replication} study when repeating a run beats trying a new one in noisy computer simulations.
\citet{chen2000ocba} split a fixed number of simulation runs among a set of options so as to pick the best one.
\citet{dai2023replicable} describe the trade-off our agents face, under large noise whose size varies across
conditions. One can test more new conditions with fewer repeats each, or fewer conditions with more repeats. Their
method sets the number of repeats for each chosen condition from the data so far. Our scripted rules are simpler: each
re-tests its current best a fixed number of times per round. A rule that sets its re-testing from the data would be a
stronger yardstick. In bandit terms, a Wellocate campaign is close to best-arm identification on a fixed budget. The
task is to find the best option with a set number of noisy tries, here made four at a time, and the score is the true
value of the final pick. In Wellocate, candidates that share letters have related values, so one
reading says something about other candidates, and the agent may nominate a candidate it never measured.

\citet{wu2016adaptation} measured the fitness of nearly every variant at four sites of the protein GB1 and mapped its
fitness peaks. Our GB1 track is built from their data. \citet{wu2019mlde} used the same landscape to test
machine-learning-guided directed evolution, which picks the next variants to measure from a model fitted to earlier
readings. Our model-guided rule is a simple version of that idea. Our grid uses the NK model \citep{kauffman1989nk}, a
classic generated landscape. Other generated landscapes come from physical models: \citet{lobkovsky2011predictability}
build them from a protein-folding model and ask how predictable evolution is on them.

We find four benchmark environments closest to our work. Wellocate combines four features: landscapes generated by program,
measurement noise whose level we set, down to zero, re-testing as a choice that uses budget, and one scored final pick.
BoxingGym \citep{gandhi2025boxinggym} is the nearest. Its ten environments are probabilistic models of fixed real-world
domains, and the agent's goal is to learn the model. Experiments are scored by expected information gain, how much
each is expected to teach about the model. The study does not vary the noise level or test a noise-free setting, and
it does not study repeating an experiment as a choice. \citet{park2026exploration} also generate their environments by
program and aim to isolate specific behaviours of agents. Their observations carry no noise, so a second look returns
what the first one showed. \citet{wang2026allocbench} have agents allocate a fixed budget, under task descriptions that
range from abstract to concrete, and their outcomes carry no noise. \citet{huai2026agenthpo} study agents that decide
from accumulated experimental results, on 30 real machine-learning tasks with no set noise level. Each of these
lacks at least one of the four features, and each keeps its noise level fixed.
